% General use
\newcommand{\RR}{\mathbb R}
\newcommand{\CC}{\mathbb C}
\newcommand{\DD}{\mathbb D}
\newcommand{\ZZ}{\mathbb Z}
\newcommand{\QQ}{\mathbb Q}
\newcommand{\PP}{\mathbb P}
\newcommand{\EE}{\mathbb E}
\newcommand{\HH}{\mathbb H}
\newcommand{\NN}{\mathbb N}

\newcommand{\comma}{\mathbin ,}
\newcommand{\conn}{\mathbin \#}
\newcommand{\sltwo}{{{\mathfrak{sl}}_2}}
\renewcommand{\sl}{\mathfrak{sl}}
\newcommand{\gl}{\mathfrak{gl}}
\newcommand{\fg}{{\mathfrak g}}
\newcommand{\co}{\colon\thinspace} % XXX: spacing is too large in subscript
\newcommand{\eps}{\varepsilon}
\newcommand{\abs}[1]{{\lvert #1 \rvert}}
\newcommand{\norm}[1]{{\lVert #1 \rVert}}
\newcommand{\OneHalf}{{\textstyle\frac{1}{2}}}

% Synonyms for commands I never remember
\newcommand{\isom}{\cong}
\newcommand{\superset}{\supset}
\newcommand{\bigcircle}{\bigcirc}
\newcommand{\contains}{\ni}
\newcommand{\tensor}{\otimes}
\newcommand{\bdy}{\partial}

% Stupid overloading.
\newcommand{\lbracket}{[}
\newcommand{\rbracket}{]}

% Various operators.
\DeclareMathOperator{\ad}{ad}
\DeclareMathOperator{\Ad}{Ad}
\DeclareMathOperator{\End}{End}
\DeclareMathOperator{\sign}{sign}
\DeclareMathOperator{\Sym}{Sym}
\DeclareMathOperator{\tr}{tr}
\DeclareMathOperator{\Hom}{Hom}
\DeclareMathOperator{\vol}{vol}
\DeclareMathOperator{\rank}{rank}
\DeclareMathOperator{\im}{im}
\DeclareMathOperator*{\esssup}{ess\,sup}
\DeclareMathOperator{\Image}{Im}

% Linear groups
\DeclareMathOperator{\ISO}{\mathit{ISO}}
\DeclareMathOperator{\SO}{\mathit{SO}}
\DeclareMathOperator{\GL}{\mathit{GL}}
\DeclareMathOperator{\SL}{\mathit{SL}}
\DeclareMathOperator{\PSL}{\mathit{PSL}}

% Special knots
\newcommand{\unknot}{\bigcircle}

% Theorems
\theoremstyle{plain}
%\numberwithin{equation}{section}
\newtheorem{proposition}[equation]{Proposition}
\newtheorem{lemma}[equation]{Lemma}
\newtheorem{corollary}[equation]{Corollary}
\newtheorem{claim}[equation]{Claim}
\newtheorem{conjecture}[equation]{Conjecture}
\newtheorem{observation}[equation]{Observation}
\newtheorem{theorem}{Theorem}
\newtheorem{citethm}[equation]{Theorem}

\theoremstyle{definition}
\newtheorem{definition}[equation]{Definition}
\newtheorem{convention}[equation]{Convention}
\newtheorem{exercise}[equation]{Exercise}
\newtheorem{question}[equation]{Question}
\newtheorem{problem}[equation]{Problem}

\theoremstyle{remark}
\newtheorem{example}[equation]{Example}
%\newtheorem{hint}[equation]{Hint}
\newtheorem{remark}[equation]{Remark}
\newtheorem{note}[equation]{Note}
%\newtheorem{apology}[equation]{Apology}
\newtheorem{warning}[equation]{Warning}

% Hyphenation.
\hyphenation{Thurs-ton}

\newcommand{\Verts}{\mathcal{V}}
\newcommand{\Edges}{\mathcal{E}}
\newcommand{\OrEdges}{\Edges^+}
\newcommand{\Lengths}{\mathcal{L}}
\newcommand{\Widths}{\mathcal{W}}
\newcommand{\Pos}{\mathcal{P}}

\newcommand{\wt}[1]{\widetilde{#1}}

\DeclareMathOperator{\EL}{EL} % extremal length
\DeclareMathOperator{\TT}{TT} % train track
\DeclareMathOperator{\SF}{SF} % stretch factor
\newcommand{\ASF}{\overline{\SF}} % asymptotic stretch factor
\DeclareMathOperator{\Fill}{Fill} % filling function
\DeclareMathOperator{\Lip}{Lip} % Lipshitz energy
\DeclareMathOperator{\Dir}{Dir} % Dirichlet energy
\DeclareMathOperator{\Emb}{Emb} % Embedding energy
\DeclareMathOperator{\Mod}{Mod} % Modulus of annulus


%%% Local Variables: 
%%% mode: latex
%%% TeX-master: t
%%% End: 
